\begin{pastexam}[label={px:ntu-109-dsaa-12}]{臺大資工所 109 學年度 第 12 題}
（問答）

(15 points) There are $n$ stations along a coastal railway ($n>0$). You're planning to select some
of them to open cafes. Three arrays $S$, $L$ and $R$ have been given, including
\begin{itemize}
\item $S=(s_1,s_2,\ldots,s_n)$: the list of the stations from $s_1$ (first) to $s_n$ (last).
\item $L=(l_1,l_2,\ldots,l_n)$: the locations of the stations, where $l_i$ is the distance of
$s_i$ from the first station $s_1$. So $l_1=0$ and $l_n$ is the length of the railway.
$l_1<l_2<\ldots<l_n$.
\item $R=(r_1,r_2,\ldots,r_n)$: the revenues of the cafes, where $r_i$ ($>0$) is the revenue for
opening a cafe in $s_i$.
\end{itemize}
The only one constraint in your plan is that the distance of any pair of your selected stations
should be longer than a given threshold $T$. If $s_i$ and $s_j$ ($i\ne j$) are selected, the
total revenue would be $l_i+l_j$. Different selection leads to different total revenue. Given
$S$, $L$, $R$ and $T$, your goal is to pick up a subset of the stations to maximize the total
revenue under the constraint. Suppose that $f(n)$ returns the maximum total revenue for the
cafes you select from the first $n$ stations. Please answer the following questions.
\textbf{No code is required (code will not be graded).}
\begin{enumerate}[label=(\alph*)]
\item (9 points) Give an $O(n^2)$ solution by defining a recurrence formula for $f(n)$. Clearly
explain the meaning of your formula and why it can be computed in $O(n^2)$ time.
\item (6 points) Consider the special case that the distance difference between two consecutive
stations is 1, i.e., $l_{i+1}-l_i=1$ ($1\le i\le n-1$). Give an $O(n)$ solution by defining a
recurrence formula for $f(n)$. Clearly explain the meaning of your formula and why it can be
computed in $O(n)$ time.
\end{enumerate}
\end{pastexam}

\begin{pastsoln}
\textbf{原題的筆誤}。 題目說選了 $s_i$、$s_j$ 時「total revenue would be $l_i + l_j$」，但 $l_i$ 是位置，
前面才定義 $r_i$ 是在 $s_i$ 開店的收益；照位置算收益沒有道理（而且 $R$ 就完全用不到了）。 這裡採用
「總收益是選到的車站的 $r$ 加總」，也就是選 $s_i$、$s_j$ 時是 $r_i + r_j$。 另外 $T$ 當成非負的數。

答案：
\begin{itemize}
\item (a) $f(0) = 0$，$f(i) = \max\bigl(f(i-1),\ r_i + f(p(i))\bigr)$，$p(i)$ 是 $l_i - l_j > T$ 的最大
$j < i$（沒有就是 $0$）；逐一往回找 $p(i)$，$O(n^2)$。
\item (b) $f(i) = \max\bigl(f(i-1),\ r_i + f(\max(0, i - d))\bigr)$，$d = \lfloor T \rfloor + 1$；$O(n)$。
\end{itemize}

\textbf{(a) 子問題}（資結書第~18.1.5~節的步驟）。 令 $f(i)$ 是只能從前 $i$ 個車站 $s_1, \ldots, s_i$ 選時
最大的總收益，$f(0) = 0$（不選）；要的答案是 $f(n)$。 車站照位置排好，選出來的車站兩兩距離都 $> T$，
等價於「相鄰兩個選到的車站距離 $> T$」：位置是遞增的，相鄰的距離 $> T$，隔得更遠的距離只會更大。

\textbf{遞迴式}。 看 $s_i$ 選不選：
\begin{itemize}
\item 不選 $s_i$：剩下的問題是從前 $i - 1$ 個選，最好是 $f(i - 1)$。
\item 選 $s_i$：其他選到的車站都在 $s_i$ 前面，而且和 $s_i$ 的距離要 $> T$，也就是只能從
$\{s_j : l_i - l_j > T\}$ 裡選。 因為 $l$ 遞增，這些車站正好是前 $p(i)$ 個，其中
\[
p(i) = \max\{\, j : 1 \le j < i,\ l_i - l_j > T \,\} \quad (\text{沒有這樣的 } j \text{ 時 } p(i) = 0).
\]
前 $p(i)$ 個裡面怎麼選才合法，和原問題一樣，最好是 $f(p(i))$。 所以這種情形最好是 $r_i + f(p(i))$。
\end{itemize}
兩種取大的：
\[
f(0) = 0, \qquad f(i) = \max\bigl(f(i - 1),\ r_i + f(p(i))\bigr) \quad (1 \le i \le n).
\]
最佳子結構：選 $s_i$ 時，扣掉 $s_i$ 剩下的選法一定是前 $p(i)$ 個車站的最佳選法；如果不是，換成更好的那個，
再加上 $s_i$ 仍然合法（新選的都在前 $p(i)$ 個裡，和 $s_i$ 的距離 $> T$），總收益更大，矛盾。 不選 $s_i$
的情形同理。

\textbf{計算和時間}。 $f(i)$ 只用到 $f$ 在比 $i$ 小的位置的值，由 $i = 1$ 填到 $n$ 就好（由下而上，資結書
第~18.1.3~節）。 每個 $i$ 從 $j = i - 1$ 往回找到第一個 $l_i - l_j > T$ 的 $j$，最多看 $i - 1$ 個，$O(n)$；
算 $f(i)$ 本身 $O(1)$。 $n$ 個 $i$ 共 $O(n^2)$ 時間、$O(n)$ 空間。 要印出選了哪些車站，從 $i = n$ 往回：
$f(i) = f(i - 1)$ 就不選 $s_i$、跳到 $i - 1$，否則選 $s_i$、跳到 $p(i)$。

（也可以寫成「$g(i)$ 是最後一個選 $s_i$ 的最大收益，$g(i) = r_i + \max\bigl(0, \max_{j : l_i - l_j > T}
g(j)\bigr)$，$f(n) = \max_i g(i)$」，每個 $g(i)$ 掃過前面所有 $j$，一樣是 $O(n^2)$。）

\textbf{(b) 間距都是 $1$}。 這時 $l_i = i - 1$，$l_i - l_j = i - j$。 令 $d = \lfloor T \rfloor + 1$，是大於
$T$ 的最小整數；$i - j > T$ 若且唯若 $i - j \ge d$，也就是 $j \le i - d$。 所以 $p(i)$ 不用找，直接是
\[
p(i) = \max(0,\ i - d),
\]
遞迴式變成
\[
f(0) = 0, \qquad f(i) = \max\bigl(f(i - 1),\ r_i + f(\max(0,\ i - d))\bigr) \quad (1 \le i \le n).
\]
意思和 (a) 一樣：不選 $s_i$，或選 $s_i$ 並且前一個選到的車站至少要在 $d$ 站之前。 每個 $f(i)$ 是 $O(1)$
（$p(i)$ 用一次減法算出），$n$ 個共 $O(n)$。 例如 $T = 2$ 時 $d = 3$，選了 $s_i$ 就不能選 $s_{i-1}$、
$s_{i-2}$；$T = 0.5$ 時 $d = 1$，相鄰的車站可以同時選，全選最好。

\textbf{補充}。 一般情形其實也能 $O(n)$：$l$ 遞增，所以 $p(i)$ 隨 $i$ 不會變小，用一個指標從上一個
$p(i - 1)$ 往前推就好，指標總共只往前走 $n$ 步。 題目 (a) 只要求 $O(n^2)$，上面的逐一往回找已經足夠。
\end{pastsoln}
